\documentclass[11pt]{article}
\usepackage[a4paper,margin=28mm]{geometry}
\usepackage[T1]{fontenc}
\usepackage{lmodern,amsmath,amssymb,amsthm,mathtools,hyperref}
\newcommand{\ddc}{dd^c}
\newcommand{\MA}{\operatorname{MA}}
\newcommand{\PSH}{\operatorname{PSH}}
\newcommand{\E}{\mathcal E}
\newtheorem{theorem}{Theorem}
\newtheorem{lemma}[theorem]{Lemma}
\theoremstyle{remark}\newtheorem{remark}[theorem]{Remark}
\title{Non-$L^p$ selection and scalar centering at zero lower mass}
\author{}\date{}
\begin{document}\maketitle
\begin{abstract}
On a product of a Hermitian manifold of zero lower Monge--Amp\`ere mass
with projective space, a separated probability density can belong to
$L^1$ and to no $L^p$, $p>1$. For an explicit radial family we determine
the limit of the entire supremum-normalized greatest exponential family.
A finite first moment gives a mean-zero selected critical potential and
an explicit finite scalar error. At the exact moment boundary and beyond,
the scalar error diverges, although the normalization scalars converge
and the normalized potentials still have a finite critical limit.
The conclusions are stable when the radial density parameter and exponent
vary together, including approaches to the moment boundary from below.
The theorem concerns this separated family; arbitrary $L^1$ data and
varying singular references remain open.
\end{abstract}

Use $\ddc=i\partial\bar\partial/(2\pi)$ and ordinary wedge powers.
All potentials are genuine canonical quasi-psh functions, and extended
products are the non-pluripolar products with smooth-background corrections.
Exponential weights multiply the actual measure $q_\alpha\kappa^k$;
they are zero at its pole. Capacity convergence uses comparison sets at
common poles, without subtracting two minus-infinity values.
The original S05 background supplies greatest finite-mass solutions
$U_\lambda$ when the individual threshold $b_H(f)$ is positive, together
with
\begin{equation}\label{eq:background}
 c_\lambda=e^{\lambda\sup U_\lambda}\longrightarrow b_H(f),
 \qquad e^{\lambda U_\lambda}f\longrightarrow b_H(f)f
 \quad\hbox{in }L^1.
\end{equation}
These scalar and measure assertions are premises here; they do not assert
convergence of normalized potentials.

\section{The family and the selection statement}
Let $(X_0,h_0)$ be a connected compact Hermitian $m$-fold with
$v_-(h_0)=0$, and set $V_0=\int_{X_0}h_0^m$. On
$X=X_0\times Y$, $Y=\mathbf P^k$, $k\geq2$, put
\[
 H=h_0+\kappa,\quad N=m+k,\quad b=\binom Nm V_0,\qquad
 dV=\frac{h_0^m}{V_0}\otimes\kappa^k,\quad\int_Y\kappa^k=1.
\]
Pullbacks are implicit. Let $o=[1:0:\cdots:0]$ and, for $0<\alpha<1$, define
\[
 \psi=\log\frac{\sum_{i=1}^k|z_i|^2}{\sum_{i=0}^k|z_i|^2},
 \qquad u_\alpha=1-(1-\psi)^\alpha,\qquad
 \MA_\kappa(u_\alpha)=q_\alpha\kappa^k.
\]
The canonical value at $o$ is $-\infty$. We verify below that the last
measure is a probability, with no pole atom.

\begin{theorem}\label{thm:selection}
The product has $v_-(H)=0$. The density $f(x,y)=q_\alpha(y)$ is in $L^1(dV)$
and in no $L^p(dV)$, $p>1$, and $b_H(f)=b>0$.
The greatest family has the form
\[
 U_\lambda=\lambda^{-1}\log b+w_\lambda\circ\pi_Y.
\]
Writing $a_\lambda=\sup_Yw_\lambda$ and $v_\lambda=w_\lambda-a_\lambda$,
the entire family $v_\lambda$ converges to $u_\alpha$ in $L^1(Y)$,
in K\"ahler capacity on $Y$, and locally uniformly on $Y\setminus\{o\}$.
The pullbacks converge in $L^1(X)$ and locally uniformly away from
$X_0\times\{o\}$.

If $0<\alpha<k/(k+1)$, then
\[
 M_\alpha=\int_Yu_\alpha q_\alpha\kappa^k\in(-\infty,0),\qquad
 a_\lambda\longrightarrow-M_\alpha,\qquad
 w_\lambda\longrightarrow u_\alpha-M_\alpha.
\]
The last convergence has the same $L^1$ and local uniform topology.
The limit is the greatest separated critical subsolution with
$q_\alpha$-mean at most zero.

If $k/(k+1)\leq\alpha<1$, then $a_\lambda\to+\infty$ and
$w_\lambda\to+\infty$ locally uniformly on $Y\setminus\{o\}$.
In both cases $c_\lambda=b e^{\lambda a_\lambda}\to b$.
In particular the actual scalar error
\[
 \lambda^{-1}\log(c_\lambda/b)=a_\lambda
\]
has the stated finite limit or diverges. The supremum-normalized limit
solves the exact threshold equation with finite mass $b$ and belongs to
the zero-reference stretching envelope class.
\end{theorem}

\section{The actual radial measure and its moment}
On the affine chart write $s=\log|z|^2$,
$m(s)=e^s/(1+e^s)$ and $\rho(s)=\log(1+e^s)$. If
$L_\alpha=(\rho+u_\alpha)'$ and $Q_\alpha=L_\alpha^k$, the radial
determinant gives
\[
 0<L_\alpha<1,\quad L_\alpha'>0,\qquad
 \MA_\kappa(u_\alpha)(|z|^2\leq e^s)=Q_\alpha(s),\qquad
 q_\alpha(s)=\frac{Q_\alpha'(s)}{(m(s)^k)'}.
\]
The initial flux is zero and the final flux is one. Canonical finite-level
truncations give this measure on every annulus; their initial missing
flux tends to zero. Thus its actual non-pluripolar mass is one, and it
does not place mass at the pole. The form is positive since
\[
 \kappa+\ddc u_\alpha
 =(1-A)\kappa+A(\kappa+\ddc\psi)
   +\alpha(1-\alpha)(1-\psi)^{\alpha-2}
                   d\psi\wedge d^c\psi,\quad
 A=\alpha(1-\psi)^{\alpha-1}\leq\alpha,
\]
and $\kappa+\ddc\psi=\ddc\log|z|^2$ off the pole.
The density extends smoothly and positively across infinity, with value
$1-\alpha$.

For $t=-s\to\infty$,
\[
 u_\alpha\sim-t^\alpha,\quad
 L_\alpha\sim\alpha t^{\alpha-1},\quad
 Q_\alpha'\sim k(1-\alpha)\alpha^k t^{-1-k(1-\alpha)}.
\]
Hence the first moment is finite exactly when $(k+1)\alpha<k$.
The density is comparable to
$e^{kt}t^{-1-k(1-\alpha)}$ near the pole, whereas smooth volume has
radial factor $e^{-kt}dt$. This proves $L^1$ integrability and failure
of every $L^p$, $p>1$.

For smooth positive $g=g(y)$ the closed projective equation and its
pullback give
$c_H(g)=b/\int_Yg\kappa^k$. Uniform multiplicative approximation gives
the same formula for positive continuous $g$, in particular
$g=\min(q_\alpha,j)$. Taking $j\to\infty$ computes $b_H(f)=b$.
Bounded small-mass potentials on $X_0$, pulled back to $X$, have product
masses $\binom Nm$ times their masses times $\int_Y\kappa^k$.
Their masses tend to zero, proving $v_-(H)=0$.

\section{Identifying the greatest exponential family}
Choose smooth nondecreasing concave $\phi$, equal to $r$ for $r\leq1$
and constant in $[1,2]$ for $r\geq2$, and put $q_j=j\phi(q_\alpha/j)$.
Then
\[
 \min(q_\alpha,j)\leq q_j\leq\min(q_\alpha,2j),\qquad q_j\uparrow q_\alpha.
\]
The functions $q_j$ are smooth and strictly positive, including near $o$
where they are constant; they are bounded below by $1-\alpha$.
The bounded positive-exponent existence and
uniqueness theorem applies without BMP \cite[Theorem 4.5]{BGL}.
For the base solutions
\[
 (\kappa+\ddc w_{\lambda,j})^k
       =e^{\lambda w_{\lambda,j}}q_j\kappa^k,
\]
uniqueness makes them invariant under $U(k)$ fixing $o$.
They are smooth by the radial
argument below. Maximum comparison outside the sole pole gives
$w_{\lambda,j}\geq u_\alpha$, and increasing the density gives a
decreasing sequence in $j$.

The pullback plus $\lambda^{-1}\log b$ is an actual bounded product
solution. The greatest solution for $q_j$ dominates it, so is bounded
below, and its canonical compact supremum is finite. It is therefore
bounded; bounded uniqueness identifies the two. Explicitly, for bounded
solutions $a,d$, compare $a$ with $d-\varepsilon$. On
$\{a<d-\varepsilon\}$ their measures have ratio at most
$e^{-\lambda\varepsilon}<1$; \cite[Lemma 4.2]{BGL} applies with a constant
strict barrier below both. Interchanging the solutions and taking
$\varepsilon\downarrow0$ proves equality.
We need comparison also for unbounded subsolutions on the full product.
Let $W$ be a bounded continuous exponential solution for $q_j$, and
let $s$ be any genuine exponential subsolution for $q_j$.
Then $t=\max(s,W)$ is bounded. Truncate $s$ below
$\inf W-1$; locality on $\{s>W\}$ and the bounded maximum/contact
inequality \cite[Lemma 1.2]{GL} on its complement give
$\MA_H(t)\geq e^{\lambda t}q_j\,dV$.
If $M=\sup(t-W)>0$, put
$m_\epsilon=\inf(W-(1-\epsilon)t)$ and
$A_{\epsilon,\tau}=\{W<(1-\epsilon)t+m_\epsilon+\tau\}$.
On this set $t-W\geq M-2\epsilon\|t\|_\infty-\tau$.
Choose a fixed sufficiently small $\epsilon>0$, and then small $\tau>0$,
so that
\[
 \MA_H((1-\epsilon)t)\geq
 (1-\epsilon)^N e^{\lambda M/2}\MA_H(W),\qquad
 (1-\epsilon)^N e^{\lambda M/2}>1
\]
there. The modified Hermitian comparison theorem
\cite[Theorem 0.2]{KN} bounds its integral by
$(1+C\tau/\epsilon^N)\int_{A_{\epsilon,\tau}}\MA_H(W)$.
This last integral is positive: a nonempty strict superlevel set of a
canonical quasi-psh function relative to a continuous function has
positive volume by the local submean inequality, and $q_j$ has a
positive lower bound. Letting $\tau\downarrow0$ contradicts the strict
factor. Thus every such subsolution is at most $W$.
Every $q_\alpha$ subsolution is also a $q_j$ subsolution, so the original
greatest family is at most every separated approximation. The limiting
solution constructed next is a $q_\alpha$ subsolution and hence is
dominated by that greatest family. This proves equality.

For fixed $\lambda$ the limit $w_\lambda$ lies between $u_\alpha$ and
$\sup w_{\lambda,1}$. Its densities are dominated by
$e^{\lambda\sup w_{\lambda,1}}q_\alpha$. Their ordinary $L^1$ limit gives
the exact identities
\begin{equation}\label{eq:base}
 \MA_\kappa(w_\lambda)=e^{\lambda w_\lambda}q_\alpha\kappa^k,\qquad
 \int_Ye^{\lambda w_\lambda}q_\alpha\kappa^k=1.
\end{equation}
These are actual products: the limiting radial flux has initial value
zero. No assertion follows merely from convergence of total masses.

\begin{lemma}\label{lem:radial}
A bounded radial solution for smooth positive $q_j$ is smooth on
$\mathbf P^k$. The limiting $w_\lambda$ is smooth off $o$ and satisfies
\[
 w_\lambda'(s)=Q_\lambda(s)^{1/k}-m(s),\quad
 Q_\lambda(s)=\int_{-\infty}^s
                   e^{\lambda w_\lambda(r)}q_\alpha(r)\,d(m(r)^k).
\]
\end{lemma}
\begin{proof}
For a bounded radial potential $\Phi=w+\rho$ is convex and increasing.
Symmetric convolution in $s$ preserves these properties and decreases
to $\Phi$ as its scale decreases. The local bounded monotone
Bedford--Taylor theorem \cite{BT} passes the smooth identity
$d(\Phi'^k)$ to the limit. Boundedness sets its initial flux to zero.
The positive density makes the flux continuous and strictly between zero
and one at finite $s$, so the derivative identity bootstraps there.

At the origin set $y=e^s$, $W(y)=w(s)$ and $Q=y^k A(y)$. Then
\[
 A(y)=k\int_0^1e^{\lambda W(yt)}q_j(yt)
                  \frac{t^{k-1}}{(1+yt)^{k+1}}\,dt,\qquad
 W'(y)=A(y)^{1/k}-\frac1{1+y}.
\]
Boundedness first gives $Q=O(y^k)$ and $w'=O(y)$, hence a continuous
Lipschitz extension of $W$. The equations and $A(0)>0$ then bootstrap
it to smoothness.

At infinity instead take $y=e^{-s}$ and $1-Q=yA(y)$. With $q_\infty$
the smooth density in this variable,
\[
 A(y)=k\int_0^1\frac{e^{\lambda W(yt)}q_\infty(yt)}
                              {(1+yt)^{k+1}}\,dt,\qquad
 W'(y)=-\frac1{1+y}+A(y)R(yA(y)),
\]
where $R(z)=[1-(1-z)^{1/k}]/z$, $R(0)=1/k$, is smooth near zero.
The tail estimate $|w'|\leq Ce^{-s}$ gives the continuous extension
and these equations bootstrap it. The variable
$y=|z_0|^2/\sum_{i=1}^k|z_i|^2$ is smooth near the hyperplane.
For the decreasing limit, dominated densities give uniform flux
convergence; integrate the derivative identity on any finite interval.
The same bootstrap applies off $o$, including infinity.
\end{proof}

\section{Convergence and the centering dichotomy}
By \eqref{eq:background}, for small $\lambda$ the actual base measures obey
\[
 \MA_\kappa(v_\lambda)
  =\frac{c_\lambda}{b}e^{\lambda v_\lambda}q_\alpha\kappa^k
  \leq2q_\alpha\kappa^k,
\]
and converge in total variation to $q_\alpha\kappa^k$.
Both $v_\lambda$ and $u_\alpha$ are supremum-normalized and full mass.
Dinew's full-mass stability theorem \cite[Theorem 1.3]{Dinew} therefore
gives $L^1$ and K\"ahler capacity convergence along every sequence
$\lambda\to0$. The common dominating measure is non-pluripolar.
This closed-class theorem is applied only on $Y$.

Uniform flux convergence gives uniform derivative convergence on finite
$s$ intervals. The $L^1$ limit fixes the integration constants there.
Near infinity, common domination and smoothness of $q_\alpha$ give
$1-Q_\lambda(s)\leq Ce^{-s}$ and $|v_\lambda'(s)|\leq Ce^{-s}$.
Integration supplies local uniform convergence including infinity.

Suppose first $M_\alpha>-\infty$. From $u_\alpha\leq w_\lambda$ and
\eqref{eq:base}, $e^r\geq1+r$ gives
$\int_Yw_\lambda q_\alpha\kappa^k\leq0$.
On a fixed annulus $A$ away from $o$, the uniform smooth-volume $L^1$
bound for $v_\lambda$ and boundedness of $q_\alpha$ imply
\[
 \int_Aw_\lambda q_\alpha\kappa^k
       \geq a_\lambda\int_Aq_\alpha\kappa^k-C.
\]
On its complement $w_\lambda\geq u_\alpha$.
Thus $0\leq a_\lambda\leq C_*$ uniformly. Any subsequential limit $a$
of these constants gives $w_\lambda\to u_\alpha+a$ off the pole.
In the identity
\[
 \int_Y\frac{e^{\lambda w_\lambda}-1}{\lambda}
                       q_\alpha\kappa^k=0,
\]
the negative part is dominated by $-u_\alpha$, and the positive part
by $C_*e^{C_*}$. Dominated convergence yields $M_\alpha+a=0$.
It identifies every subsequence and gives the asserted global $L^1$
convergence by the same integrable bounds.
No unit-mass Jensen inequality is used for the approximants $q_j$;
their masses differ from one.

If $M_\alpha=-\infty$, \eqref{eq:base} instead gives
\[
 \int_Y\frac{1-e^{\lambda v_\lambda}}{\lambda}
                q_\alpha\kappa^k
   =\frac{1-e^{-\lambda a_\lambda}}{\lambda}\leq a_\lambda.
\]
The integrands are nonnegative and tend pointwise off $o$ to
$-u_\alpha$. Fatou's lemma forces $a_\lambda\to+\infty$.
The locally uniform limit of $v_\lambda$ then proves the divergence
of $w_\lambda$ off the pole. Meanwhile $\lambda a_\lambda\to0$ by
\eqref{eq:background}. This proves the dichotomy without discarding
the error term.

In the finite-moment regime the constant is explicitly
\[
 -M_\alpha=\alpha\int_1^\infty x^{\alpha-1}
 \left[\alpha x^{\alpha-1}
       +(1-\alpha x^{\alpha-1})e^{1-x}\right]^k\,dx.
\]
This follows by integrating $u_\alpha\,dQ_\alpha$ by parts with
$x=1-\psi$; the pole boundary term vanishes because
$(k+1)\alpha-k<0$.

\section{Joint variation of the density and exponent}

\begin{theorem}\label{thm:joint}
In the fixed product above, let $\alpha_j\to\alpha\in(0,1)$ and
$\lambda_j\downarrow0$, without any restriction on their relative rates.
Then $q_{\alpha_j}\to q_\alpha$ in $L^1$, and
$v_{\lambda_j,\alpha_j}\to u_\alpha$ in $L^1(Y)$, in $\kappa$ capacity
and locally uniformly off $o$, including infinity. The product pullbacks
converge in $L^1(X)$ and locally uniformly off the pole fiber.
Their actual scalars satisfy $c_j/b\to1$.

If $\alpha<k/(k+1)$, then $a_{\lambda_j,\alpha_j}\to-M_\alpha$ and
$w_{\lambda_j,\alpha_j}\to u_\alpha-M_\alpha$ in $L^1$ and locally
uniformly off the pole. If $\alpha\geq k/(k+1)$, then
$a_{\lambda_j,\alpha_j}\to+\infty$, $\lambda_j a_{\lambda_j,\alpha_j}\to0$
and $w_{\lambda_j,\alpha_j}\to+\infty$ locally uniformly off the pole.
This includes approaching the critical parameter from below.
The reference is zero throughout; no moving singular model is asserted.
\end{theorem}
\begin{proof}
Choose a compact parameter interval $[\alpha_-,\alpha_+]\subset(0,1)$
containing $\alpha$ and every late $\alpha_j$. The explicit radial formulas,
uniformly on this interval, give for $t=-s$ large
\[
 q_\delta\leq C e^{kt}(1+t)^{-1-k(1-\alpha_+)}.
\]
Indeed $L_\delta\asymp\delta t^{\delta-1}$ and
$L_\delta'\asymp\delta(1-\delta)t^{\delta-2}$ uniformly, while
$(m^k)'\asymp e^{-kt}$. On the compact complement the densities are
jointly smooth, including infinity. Thus one nonnegative $L^1$ density
$g$ dominates every $q_\delta$ in this interval. Its finite measure
is non-pluripolar. Dominated convergence gives the asserted density limit.

Normalized $\kappa$-psh functions satisfy the canonical compactness bound
$\int(-v)\kappa^k\leq C_k$. Consequently
$\kappa^k(\{v<-R\})\leq C_k/R$. Absolute continuity of the integral of
$g$ gives $\sup_v\int_{\{v<-R\}}g\kappa^k\to0$.
For $R=\lambda^{-1/2}$, therefore,
\[
 \int(1-e^{\lambda v})q_\delta\kappa^k
 \leq\lambda R+\int_{\{v<-R\}}g\kappa^k\longrightarrow0
\]
uniformly in $\delta$ and all normalized $\kappa$-psh $v$.
This uses no uniform first moment against $g$.
The actual full-mass equation says
$1=e^{\lambda_j a_j}\int e^{\lambda_j v_j}q_{\alpha_j}\kappa^k$.
It proves $e^{\lambda_j a_j}\to1$, hence $\lambda_j a_j\to0$ and
$c_j/b\to1$. The equation measures converge in total variation to
$q_\alpha\kappa^k$ and are bounded by $2g\kappa^k$ for every late $j$.
Theorem 1.3 of \cite{Dinew} applies to this single finite non-pluripolar
dominating measure and the actual supremum-zero full-mass potentials.
It gives $L^1$ and base capacity convergence.
Uniform cumulative-flux convergence, the derivative identity and
$1-Q_j(s)\leq Ce^{-s}$ near infinity give the same local uniform upgrade
as before. Pullback preserves the claimed $L^1$ and local uniform limits.

For $\alpha<k/(k+1)$ choose $\alpha_+<k/(k+1)$.
Then $(-u_{\alpha_+})g$ is integrable, since its radial power is
$\alpha_+-1-k(1-\alpha_+)<-1$.
The bound $w_j\geq u_{\alpha_j}\geq u_{\alpha_+}$, a fixed positive-mass
annulus and the normalized smooth-volume $L^1$ bound give
$0\leq a_j\leq C$, as in the fixed-parameter proof.
In the exact difference-quotient identity the negative part is
dominated by $(-u_{\alpha_+})g$ and the positive part by $Ce^C g$.
Every convergent subsequence of $a_j$ therefore has limit $-M_\alpha$.
The already proved $L^1$ convergence of $v_j$ gives that of $w_j$.

For $\alpha\geq k/(k+1)$ use instead
\[
 \int\frac{1-e^{\lambda_j v_j}}{\lambda_j}q_{\alpha_j}\kappa^k
 =\frac{1-e^{-\lambda_j a_j}}{\lambda_j}\leq a_j.
\]
Off the pole the nonnegative integrands tend to
$(-u_\alpha)q_\alpha$. Fatou's lemma forces $a_j\to+\infty$, even when
every $\alpha_j$ is below the moment boundary.
Together with $\lambda_j a_j\to0$ and local uniform convergence of
$v_j$, this proves the remaining assertions.
\end{proof}

\section{Order selection and the envelope class}
Any canonical critical subsolution on $Y$ has extended mass at most one:
each closed-background bounded truncation has mass one. If its measure
dominates $q_\alpha\kappa^k$, it has that exact full measure.
Full-mass uniqueness \cite[Theorem 1.2]{Dinew} makes it $u_\alpha+d$.
Thus $u_\alpha$ is the unique supremum-zero member, giving the normalized
selection an order characterization throughout $0<\alpha<1$.
For finite $M_\alpha$, its $q_\alpha$-mean is nonpositive exactly when
$d\leq-M_\alpha$. Hence $u_\alpha-M_\alpha$ is the greatest such member
in the separated sector.

For every $t\geq1$ choose $\alpha<\beta<1$. The function
$u_\beta-tu_\alpha$ tends to $-\infty$ at $o$ and has a finite
maximum $C_t$ elsewhere. Thus $a_t=u_\beta-C_t$ is genuine and
$a_t/t\leq u_\alpha$, canonically also at the common pole.
Its product pullback proves the all-scale stretching definition of
$\E_{\rm env}(X,H,0)$. The exact product mass is
$\binom Nm V_0=b$. No uniqueness of arbitrary nonclosed critical
potentials or varying-reference theorem is asserted.

\begin{remark}
For a concrete application use the elliptic quotient of
$\{[z:w]\in\mathbf P^3:z,w\in\mathbf C^2\setminus\{0\}\}$ by
$[z:w]\mapsto[z:e^{2\pi}w]$.
With $t=\log(\|w\|^2/\|z\|^2)$ and normalized base forms
$\alpha_0,\beta_0$ on $\mathbf P^1\times\mathbf P^1$, take
$h_0=\alpha_0+\beta_0+dt\wedge d^c t$.
The complete factor construction and smooth mass-collapse proof are in
\cite[Appendix A]{Product}; they give $v_-(h_0)=0$ and $V_0=24\pi$.
For $k=2$ the theorem applies on a concrete fivefold with
$b=240\pi$. Choosing $\alpha=1/2$ gives finite scalar centering, while
$\alpha\geq2/3$ gives divergence of the scalar error.
\end{remark}

\begin{thebibliography}{9}
\bibitem{BT} E. Bedford and B. A. Taylor,
\emph{A new capacity for plurisubharmonic functions},
Acta Math. 149 (1982), 1--40.
\bibitem{BGL} S. Boucksom, V. Guedj and C. H. Lu,
\emph{Volumes of Bott--Chern classes}, arXiv:2406.01090v1.
\bibitem{Dinew} S. Dinew,
\emph{Uniqueness and stability in $\E(X,\omega)$},
arXiv:0804.3407v1; J. Funct. Anal. 256 (2009), 2113--2122.
\bibitem{GL} V. Guedj and C. H. Lu,
\emph{Quasi-plurisubharmonic envelopes 2: bounds on Monge--Amp\`ere volumes},
arXiv:2106.04272v1.
\bibitem{KN} S. Ko\l odziej and N. C. Nguyen,
\emph{Weak solutions to the complex Monge--Amp\`ere equation on Hermitian
manifolds}, arXiv:1312.5491v2.
\bibitem{Product} \emph{Exact equations for $L^1$ densities on products
with zero lower mass}, companion construction, Appendix A;
source: \texttt{P07-zero-lower-mass-product-equations.tex}.
\end{thebibliography}
\end{document}
